\documentclass[11pt]{article}
\usepackage{explorer}
\newcommand{\ExplorerLanguage}{FR}
\newcommand{\ExplorerNumber}{01}
\newcommand{\ExplorerKind}{Exercices progressifs}
\begin{document}
\ExplorerTitle{Seconde · raisonnement}{Parité et invariants}
\ExplorerSection{Exercice 1}
On part de 7 et on ajoute ou retranche 4 à chaque coup. Peut-on atteindre 20 ? Et 23 ?\par
\vspace{5pt}
\ExplorerSection{Exercice 2}
Un produit de cinq nombres égaux à 1 ou −1 vaut −1. Montrer qu’il est impossible de rendre tous les nombres positifs en changeant deux signes à chaque coup.\par
\vspace{5pt}
\ExplorerSection{Exercice 3}
On remplace deux entiers a et b par \(a + b\). Quel invariant conserve cette opération ?\par
\vspace{5pt}
\ExplorerSection{Exercice 4}
Pourquoi deux états ayant la même parité ne sont-ils pas nécessairement reliés par les coups autorisés ? Donner un exemple.\par
\vspace{5pt}
\end{document}
